\documentclass[aspectratio=169,10pt]{beamer}

\usetheme{Madrid}
\usecolortheme{default}
\setbeamertemplate{navigation symbols}{}
\setbeamertemplate{footline}[frame number]

\usepackage[utf8]{inputenc}
\usepackage[T1]{fontenc}
\usepackage{lmodern}
\usepackage{amsmath,amssymb,mathtools}
\usepackage{booktabs}
\usepackage{tabularx}
\usepackage{array}
\usepackage{siunitx}
\usepackage{xcolor}
\usepackage{graphicx}
\usepackage{adjustbox}
\usepackage{hyperref}
\usepackage{ragged2e}

\definecolor{UChicagoMaroon}{RGB}{128,0,0}
\definecolor{SoftGray}{RGB}{245,245,245}
\definecolor{PaleMaroon}{RGB}{250,238,238}
\definecolor{DarkText}{RGB}{35,35,35}
\definecolor{MutedBlue}{RGB}{40,80,130}

\setbeamercolor{structure}{fg=UChicagoMaroon}
\setbeamercolor{block title}{fg=white,bg=UChicagoMaroon}
\setbeamercolor{block body}{bg=SoftGray,fg=black}
\setbeamercolor{alerted text}{fg=UChicagoMaroon}
\setbeamerfont{frametitle}{size=\large,series=\bfseries}
\setbeamerfont{block title}{size=\small,series=\bfseries}
\setbeamerfont{block body}{size=\small}

\newcommand{\TOF}{\operatorname{TOF}}
\newcommand{\LED}{\operatorname{LED}}
\newcommand{\CFD}{\operatorname{CFD}}
\newcommand{\CTR}{\operatorname{CTR}}
\newcommand{\ps}{\,\mathrm{ps}}
\newcommand{\ns}{\,\mathrm{ns}}

% Figure paths are resolved relative to \figroot.
% Use "." when main.tex and figures/ are in the same directory.
% If needed, change ONLY this line, e.g.:
%   \renewcommand{\figroot}{..}
%   \renewcommand{\figroot}{presentation}
\newcommand{\figroot}{.}

% Figure wrapper: keeps compilation robust when a plot is missing.
\newcommand{\smartimage}[4][\linewidth]{%
  \edef\resolvedfigure{\figroot/#3}%
  \IfFileExists{\resolvedfigure}{%
    \includegraphics[width=#1,height=#2,keepaspectratio]{\resolvedfigure}%
  }{%
    \fbox{\parbox[c][#2][c]{#1}{\centering\color{DarkText}\small
    \textbf{Missing figure}\\[0.3em]
    \texttt{\detokenize\expandafter{\resolvedfigure}}\\[0.3em]
    #4}}
  }%
}

\newcommand{\takeaway}[1]{%
  \vspace{0.25em}
  {\footnotesize\color{UChicagoMaroon}\textbf{Takeaway:} #1\par}
}

\newcommand{\mainmsg}[1]{%
  \vspace{0.2em}
  \begin{center}
  \color{UChicagoMaroon}\bfseries\large #1
  \end{center}
}

\newcommand{\repo}{\href{https://github.com/Alessandro-F-GH/tof-pet-timing-uchicago-ONAOSI-prj3.git}{\texttt{tof-pet-timing-uchicago-ONAOSI-prj3}}}

\title[CTR timing analysis]{\LARGE Performance characterization of PET detectors and systems}
\subtitle{\large Oscilloscope baseline, Pico-TDC performance, and waveform-ML correction}
\author{\Large Alessandro Falcetta}
\institute{\large University of Chicago\\[0.45em]\large \textbf{Tutor:} Hee-Jong Kim\\\large \textbf{PI:} Chien-Min Kao}
\date{\large August 2026}

\begin{document}

% ---------------------------------------------------------
\begin{frame}[plain]
  \titlepage
\end{frame}

% ---------------------------------------------------------
\begin{frame}{Why timing matters in TOF-PET}
\small

\begin{columns}[T,totalwidth=\textwidth]

\begin{column}{0.48\textwidth}

\begin{block}{TOF-PET imaging}
A positron annihilation produces two photons emitted in opposite directions.
Two detectors measure the corresponding signals and estimate the
\textbf{difference in photon arrival time (TOF)}.
\end{block}

\vspace{0.25em}

\begin{alertblock}{Why better timing helps}
The arrival-time difference constrains the annihilation position along the line connecting the two detectors:

\[
\Delta x = \frac{c\,\Delta t}{2}.
\]

\vspace{0.15em}
A more precise estimate of $\Delta t$ means a more precise localization of the event.
\end{alertblock}

\end{column}


\begin{column}{0.48\textwidth}
\centering

\smartimage[0.94\linewidth]{0.57\textheight}
{figures/TOF.png}
{TOF-PET principle: the difference in photon arrival times localizes the annihilation event along the line of response.}

\vspace{0.3em}

\begin{block}{Project question}
\centering
What is the best timing performance achievable,
and how can we approach it with a practical and scalable method?
\end{block}

\end{column}

\end{columns}

\end{frame}
% ---------------------------------------------------------
\begin{frame}{Experimental setup}
\small
\begin{columns}[T,totalwidth=\textwidth]
\begin{column}{0.44\textwidth}
\centering
\smartimage[0.96\linewidth]{0.84\textheight}{figures/Board-img.png}{Detector modules and front-end boards.}
\end{column}

\begin{column}{0.53\textwidth}
\begin{block}{Detector system}
\begin{itemize}
  \item Two opposing LYSO--SiPM detector modules.
  \item LYSO crystal: $2\times2\times3~\mathrm{mm^3}$.
  \item BC SiPM: $4\times4~\mathrm{mm^2}$.
  \item FBK-HF front-end / amplifier.
  \item Positron-emitting source centered between detectors.
\end{itemize}
\end{block}

\begin{block}{Analog outputs}
Each detector side provides:
\begin{itemize}
  \item \textbf{energy} channel;
  \item \textbf{timing} channel.
\end{itemize}
\end{block}

\end{column}
\end{columns}
\end{frame}


% ---------------------------------------------------------
\begin{frame}{Two ways to acquire the detector signal}
\small

\begin{columns}[T,totalwidth=\textwidth]

% ---------------------------------------------------------
\begin{column}{0.45\textwidth}

\begin{block}{Oscilloscope: Full waveform}

The voltage signal is sampled as a function of time.

\begin{itemize}
    \item complete signal shape available offline;
    \item flexible timing analysis;

\end{itemize}
\end{block}

\vspace{0.15em}

\centering
\smartimage[0.94\linewidth]{0.34\textheight}
{figures/waveform.png}
{Example of sampled oscilloscope waveforms.}

\end{column}

\begin{column}{0.45\textwidth}

\begin{block}{Pico-TDC: Threshold-based}

Records the times at which the signal crosses a fixed threshold.

\begin{itemize}
    \item leading and trailing crossing times;
    \item Time-over-Threshold (ToT);
\end{itemize}
\end{block}

\vspace{0.15em}

\centering
\smartimage[0.82\linewidth]{0.34\textheight}
{figures/txt_list.png}
{Example of Pico-TDC timestamp data.}

\end{column}

\end{columns}

\vspace{0.25em}

\mainmsg{
Same detector signals, two different levels of information:
full waveforms versus a small set of timestamps.
}

\end{frame}

% ---------------------------------------------------------
\begin{frame}{Why use the Pico-TDC?}
\small

\begin{columns}[T,totalwidth=\textwidth]

\begin{column}{0.53\textwidth}
\centering

\smartimage[0.92\linewidth]{0.80\textheight}
{figures/Pico-TDC-img.png}
{Pico-TDC board photograph.}

\end{column}

\begin{column}{0.44\textwidth}

\begin{block}{Practical motivation}
\begin{itemize}
    \item \textbf{Lower cost} than high-bandwidth waveform digitization.
    \item \textbf{More channels} available in a compact system.
    \item Much \textbf{smaller data volume }per event.
  
\end{itemize}
\end{block}

\vspace{0.3em}

\begin{alertblock}{Main question}
\centering
Can this readout preserve the timing performance of full waveform acquisition?
\end{alertblock}

\end{column}

\end{columns}

\end{frame}
% ---------------------------------------------------------

\begin{frame}{Experimental data}
\small

\begin{columns}[T,totalwidth=\textwidth]

% ---------------------------------------------------------
\begin{column}{0.45\textwidth}
\begin{block}{Oscilloscope}
\begin{itemize}
    \item \textbf{Bias voltage:} 45--49 V, 1 V steps
    \item \textbf{Acquisition:} 100 ns waveforms at 80 GS/s
 
\end{itemize}
\end{block}

\begin{table}
\centering
\scriptsize
\setlength{\tabcolsep}{4pt}
\renewcommand{\arraystretch}{0.95}

\begin{tabular}{rrr}
\toprule
\textbf{Bias [V]} &
\textbf{Raw events} &
\textbf{Photopeak events} \\
\midrule
45 & 20\,135 & 6\,763 \\
46 & 21\,510 & 8\,037 \\
47 & 20\,073 & 8\,482 \\
48 & 20\,264 & 7\,049 \\
49 & 25\,054 & 7\,078 \\
\bottomrule
\end{tabular}

\end{table}

\end{column}

% ---------------------------------------------------------
\begin{column}{0.45\textwidth}
\begin{block}{Pico-TDC}
\begin{itemize}
    \item \textbf{Bias voltage:} 42--49 V, 1 V steps
    \item \textbf{Timing threshold:} 20--200 mV
    \item \textbf{Acquisition:} trigger-matching mode

\end{itemize}
\end{block}

\begin{table}
\centering
\scriptsize
\setlength{\tabcolsep}{4pt}
\renewcommand{\arraystretch}{0.95}

\begin{tabular}{rrrr}
\toprule
\textbf{Bias [V]} &
\textbf{$T_{\mathrm{th}}$ [mV]} &
\textbf{Raw events} &
\textbf{Photopeak events} \\
\midrule
42 & 40 & 79\,127 & 11\,514 \\
43 & 40 & 84\,586 & 11\,470 \\
44 & 40 & 90\,229 & 11\,554 \\
45 & 40 & 96\,008 & 11\,552 \\
45 & 40 & 96\,001 & 11\,658 \\
46 & 20 & 49\,348 & 5\,793 \\
46 & 30 & 49\,688 & 5\,909 \\
46 & 40 & 100\,916 & 11\,847 \\
46 & 50 & 50\,245 & 5\,931 \\
46 & 60 & 50\,271 & 5\,800 \\
46 & 70 & 50\,410 & 5\,981 \\
47 & 40 & 104\,916 & 11\,881 \\
48 & 40 & 107\,613 & 12\,182 \\
49 & 40 & 111\,406 & 12\,434 \\
\bottomrule
\end{tabular}
\end{table}

\end{column}

\end{columns}


\end{frame}
% ----------------------------------------------------------
\begin{frame}{Standard methods for oscilloscope waveforms}
\small
\centering
\smartimage[0.78\linewidth]{0.55\textheight}{figures/led_cfd_waveform_schematic.pdf}{LED and CFD crossing definition.}

\vspace{0.25em}
\begin{columns}[T,totalwidth=\textwidth]
\begin{column}{0.48\textwidth}
\begin{block}{Leading-edge discriminator}
\centering
$\Delta t_{\LED}=t_{\LED,1}-t_{\LED,2}$
\RaggedRight
Fixed voltage crossing: simple and robust, but sensitive to amplitude-dependent time walk.
\end{block}
\end{column}

\begin{column}{0.48\textwidth}
\begin{block}{Constant-fraction discriminator}
\centering
$\Delta t_{\CFD}=t_{\CFD,1}-t_{\CFD,2}$
\RaggedRight
Fractional crossing relative to pulse amplitude: reduces time walk, but still compresses the waveform to one time marker.
\end{block}
\end{column}
\end{columns}
\end{frame}
%---------------------------------------------------------

% ---------------------------------------------------------
\begin{frame}{CTR extraction}
\small
\begin{columns}[T,totalwidth=\textwidth]
\begin{column}{0.62\textwidth}
\centering
\vspace{0.5cm}
\smartimage[0.98\linewidth]{0.71\textheight}{figures/timing_fit.png}{Representative Gaussian timing fit.}
\end{column}

\begin{column}{0.36\textwidth}
\begin{block}{Timing difference}
For each selected coincidence event:
\[
\Delta t = t_1-t_2 .
\]
The selected events form the coincidence time distribution.
\end{block}

\begin{block}{CTR definition}
\[
\CTR = \mathrm{FWHM}=2\sqrt{2\ln2}\,\sigma \simeq 2.35\,\sigma .
\]
Lower CTR means a narrower coincidence-time distribution and better timing precision.
\end{block}
\end{column}
\end{columns}
\mainmsg{The same Gaussian-fit CTR definition is used for every timing method.}
\end{frame}

% ---------------------------------------------------------
\begin{frame}{Pico-TDC analysis: from hits to coincidences}
\small
\begin{columns}[T,totalwidth=\textwidth]
\begin{column}{0.6\textwidth}
\centering
\smartimage[0.98\linewidth]{0.41\textheight}{figures/peak_selection_ch5.png}{ToT photopeak selection.}
\vspace{0.25em}
\smartimage[0.98\linewidth]{0.44\textheight}{figures/alignment_plot_ch5.png}{Alignment and mismatch rejection.}
\end{column}

\begin{column}{0.34\textwidth}
\begin{block}{Pipeline}
\begin{enumerate}
    \item Raw timestamp lists
    \item Leading/trailing pairs
    \item ToT-based photopeak selection
    \item Timing-hit matching
    \item Mismatch rejection
    \item Gaussian CTR fit
\end{enumerate}
\end{block}

\begin{alertblock}{Personal contribution}
Causal signal matching and filtering strategy to reject mismatched or noisy events while keeping lower timing thresholds usable.
\end{alertblock}
\end{column}
\end{columns}
\end{frame}

% ---------------------------------------------------------
\begin{frame}{Pico-TDC threshold scan}
\small
\begin{columns}[T,totalwidth=\textwidth]
\begin{column}{0.58\textwidth}
\centering
\smartimage[0.82\linewidth]{0.72\textheight}{figures/CTR_vs_th.png}{CTR versus timing threshold.}
\end{column}

\begin{column}{0.39\textwidth}
\begin{block}{Threshold trade-off}
\begin{itemize}
  \item \textbf{Low threshold}: earlier pick-off, but more noise/spurious hits.
  \item \textbf{High threshold}: cleaner crossing, but more time walk and possible efficiency loss.
\end{itemize}
\end{block}

\begin{alertblock}{Analysis result}
The threshold can affect significantly the performance and should be optimized.
\end{alertblock}
\end{column}
\end{columns}
\end{frame}

% ---------------------------------------------------------
\begin{frame}{Pico-TDC versus oscilloscope reference}
\small
\begin{columns}[T,totalwidth=\textwidth]
\begin{column}{0.64\textwidth}
\centering
\vspace{0.3cm}
\smartimage[0.92\linewidth]{0.75\textheight}{figures/CTR_vs_V.png}{Pico-TDC versus oscilloscope CTR comparison.}
\end{column}

\begin{column}{0.32\textwidth}
\vspace{0.5cm}
\begin{block}{Result}
Using the timing channels, Pico-TDC can produce \textbf{CTR compatible with the oscilloscope reference}.
\end{block}

\begin{alertblock}{Takeaway}
    Pico-TDC is feasible for practical timing readout; the oscilloscope remains the baseline for studying the information limit.
\end{alertblock}

\end{column}
\end{columns}
\end{frame}

% ---------------------------------------------------------
\begin{frame}{Second question: can we do better?}

\small
\begin{columns}[T,totalwidth=\textwidth]
\begin{column}{0.36\textwidth}
\begin{block}{Idea}
The signal shape contains information about how the event was measured.
\end{block}

\begin{block}{Main question}
Can a model learn a waveform- dependent timing correction, rather than a black-box TOF prediction?
\end{block}

\begin{alertblock}{Risk with Machine Learning}
A very flexible model can make the CTR distribution narrower without necessarily learning the physical timing correction we want.
\end{alertblock}
\end{column}

\begin{column}{0.58\textwidth}
\centering

\smartimage[0.78\linewidth]{0.65\textheight}
{figures/distr_targ.png}
{Example from previous waveform-ML work.}

\vspace{0.15em}

{
Source: X.~Feng, A.~Muhashi, Y.~Onishi, R.~Ota, and H.~Liu,
``Transformer-CNN hybrid network for improving PET time of flight prediction,''
\textit{Physics in Medicine \& Biology}, 2024.
}
\end{column}
\end{columns}
\end{frame}

% ---------------------------------------------------------
\begin{frame}{What should the model learn?}
\small

% ---------------------------------------------------------
\begin{block}{Start from standard leading-edge timing}
For each coincidence event, the measured time difference is:  $\Delta t_{\LED}=t_{\LED,1}-t_{\LED,2}$.

\end{block}

\vspace{-0.15em}

% ---------------------------------------------------------
\begin{alertblock}{My modeling proposal}
\large
\[
\boxed{
\Delta t_{\LED}
=
\underbrace{\TOF}_{\text{true time of flight}}
+
\underbrace{C_{1,2}}_{\text{calibration bias}}
+
\underbrace{\delta(s_1,s_2)}_{\text{waveform-dependent error}}
+
\underbrace{\epsilon}_{\text{random noise}}
}
\]
\small

\end{alertblock}

\vspace{0.1em}

% ---------------------------------------------------------
\begin{columns}[T,totalwidth=\textwidth]

\begin{column}{0.32\textwidth}
\begin{block}{ML target}
\Large
\[
\boxed{\delta(s_1,s_2)}
\]
\small
Where $(s_1,s_2)$ is a signal pair.

\end{block}
\end{column}

\begin{column}{0.64\textwidth}
\begin{block}{Derived properties}
\begin{itemize}
    \item The physical TOF remains outside the ML model.
    \item Fixed channel offsets are handled by calibration.
    \item ML is used to correct waveform-dependent timing effects.
\end{itemize}
\end{block}
\end{column}

\end{columns}

\end{frame}

% ---------------------------------------------------------
\begin{frame}{Physically constrained correction}
\small

\begin{block}{Proposed model architecture: same model for all signals}

\large
\[
\boxed{y_\theta(s_1,s_2)=g_\theta(s_1)-g_\theta(s_2)}
\]
\small

\end{block}

\vspace{0.15em}

\begin{columns}[T,totalwidth=\textwidth]
\begin{column}{0.48\textwidth}
\begin{alertblock}{Corrected timing estimate}
\large
\[
\boxed{
\widehat{\Delta t}=\Delta t_{\LED}-\widehat C_{1,2}-y_\theta(s_1,s_2)
}
\]
\small
The final output is still a corrected arrival-time difference, so it can be evaluated with the same CTR metric.
\end{alertblock}
\end{column}

\begin{column}{0.48\textwidth}
\begin{block}{Why this is safer than generic pair-ML}
\begin{itemize}
  \item Same pulse scorer for every detector.
  \item Learns a timing-error correction, not an absolute position label.
  \item No need for multiple source positions: a single central source is sufficient.
\end{itemize}
\end{block}
\end{column}
\end{columns}

\vspace{0.25em}
\mainmsg{Constraining the model gives a translation-invariant, antisymmetric,
and more general timing correction.}
\end{frame}

% --------------------------------------------------------------
\begin{frame}{Windowing and target correction}
\small

\begin{center}
\smartimage[0.96\textwidth]{0.52\textheight}
{figures/windowing_native_grid.pdf}
{Native-grid windowing around the interpolated LED time.}
\end{center}

\vspace{-0.45em}

\begin{columns}[T,totalwidth=\textwidth]

\begin{column}{0.48\textwidth}
\begin{block}{Windowing}
LED time is interpolated, but waveform samples remain on the native grid:
\[
\delta_j=t_{\LED,j}-t_{a,j}.
\]
\end{block}
\end{column}

\begin{column}{0.48\textwidth}
\begin{alertblock}{Target}
Correct for the discrete anchor shift:
\large
\[
\boxed{
y_{\mathrm{target}}
=
\Delta t_{\LED}
-
\Delta\delta
-
TOF
}
\]
\end{alertblock}
\end{column}

\end{columns}

\end{frame}
% ---------------------------------------------------------
\begin{frame}{Analysis workflow: from detector data to result}
\small

\begin{alertblock}{One reproducible chain}
\centering
Raw data $\rightarrow$ event selection $\rightarrow$ LED/CFD timing $\rightarrow$ ML correction $\rightarrow$ blind CTR $\rightarrow$ interpretation
\end{alertblock}

\vspace{0.35em}
\begin{columns}[T,totalwidth=\textwidth]
\begin{column}{0.31\textwidth}
\begin{block}{1. Clean events}
\begin{itemize}
  \item Select photopeak events.
  \item Reject mismatched timing hits.
  \item Freeze standard LED/CFD references.
\end{itemize}
\end{block}
\end{column}

\begin{column}{0.31\textwidth}
\begin{block}{2. Compare readouts}
\begin{itemize}
  \item Full oscilloscope waveform.
  \item Pico-TDC timestamp readout.
  \item Reduced fixed-threshold representation.
\end{itemize}
\end{block}
\end{column}

\begin{column}{0.31\textwidth}
\begin{block}{3. Test fairly}
\begin{itemize}
  \item Tune on development data.
  \item Retrain selected configuration.
  \item Report one independent blind CTR.
\end{itemize}
\end{block}
\end{column}
\end{columns}

\vspace{1.45em}
\centering\normalsize Repository: \repo
\end{frame}

% ---------------------------------------------------------
\begin{frame}{Models as controlled information probes}
\small

\begin{block}{Machine learning models}
\begin{tabularx}{\linewidth}{@{}p{0.22\linewidth}p{0.32\linewidth}X@{}}
\toprule
\textbf{Model} & \textbf{Description} & \textbf{Scientific question}\\
\midrule
Linear SVR
& Weighted sum of waveform samples
& Is the correction mostly linear?\\[1em]

CNN
& Learns local waveform patterns
& Does nonlinear waveform structure contain additional timing information?\\[2.em]

Fixed-threshold SVR
& Uses only a few crossing times
& How much timing information survives a TDC-like readout?\\
\bottomrule
\end{tabularx}
\end{block}
\vspace{0.3em}

\begin{columns}[T,totalwidth=\textwidth]
\begin{column}{0.26\textwidth}
\begin{block}{Performance}
A model is useful if it reduces CTR on the independent blind set.
\end{block}
\end{column}

\begin{column}{0.66\textwidth}
\begin{alertblock}{Interpretation}
If different models learn similar event-by-event corrections,
performance is likely limited by the information in the signal,
rather than by model complexity.
\end{alertblock}
\end{column}
\end{columns}
\end{frame}

% ---------------------------------------------------------
\begin{frame}{Waveform-ML results: energy-channel timing}
\centering
\smartimage[0.95\linewidth]{0.90\textheight}{figures/results_energy_ctr_vs_voltage.pdf}{Energy-channel blind CTR versus voltage.}
\end{frame}

% ---------------------------------------------------------
\begin{frame}{Waveform-ML results: timing-channel timing}
\centering
\smartimage[0.95\linewidth]{0.90\textheight}{figures/results_timing_ctr_vs_voltage.pdf}{Timing-channel blind CTR versus voltage.}
\end{frame}

% ---------------------------------------------------------
\begin{frame}{Fixed-threshold timing as reduced readout}
\small
\begin{columns}[T,totalwidth=\textwidth]
\begin{column}{0.50\textwidth}
\centering
\smartimage[0.92\linewidth]{0.42\textheight}{figures/multithreshold_waveform_crossings.pdf}{Interpolated crossing times at fixed thresholds.}
\vspace{0.25em}
\smartimage[0.92\linewidth]{0.4\textheight}{figures/results_multithreshold_ctr_vs_threshold_count_energy_to_energy.pdf}{CTR improvement versus number of thresholds.}
\end{column}

\begin{column}{0.46\textwidth}
\begin{block}{Reduced feature vector}
\[
\tau_k=t_1(V_k)-t_2(V_k), \qquad x=[\tau_{k_1},\ldots,\tau_{k_m}] .
\]
Only $m$ signal crossing times are passed to SVR.
\end{block}

\begin{block}{Result}
On energy channels, the first few thresholds recover much of the correction and can outperform standard single-time estimators.
\end{block}

\begin{alertblock}{Connection to Pico-TDC}
This tests how much information can survive a TDC-like representation, without considering the ToT information.
\end{alertblock}
\end{column}
\end{columns}
\end{frame}

% ---------------------------------------------------------
% ---------------------------------------------------------
\begin{frame}{XAI: do different models find the same correction?}
\small

\begin{columns}[T,totalwidth=\textwidth]

\begin{column}{0.58\textwidth}
\centering
\smartimage[0.98\linewidth]{0.68\textheight}
{figures/output_correlation.png}
{Correlation between event-by-event corrections from different models.}
\end{column}

\begin{column}{0.38\textwidth}

\begin{block}{Result}
SVR, CNN, and fixed-threshold SVR produce
\textbf{strongly correlated event-by-event corrections}.
\end{block}

\vspace{0.4em}

\begin{alertblock}{Interpretation}
Similar corrections across different models suggest that
signal information, not model complexity, limits performance.
\end{alertblock}

\end{column}
\end{columns}

\end{frame}

% ---------------------------------------------------------
\begin{frame}{XAI: energy-channel information}
\small
\centering
\smartimage[0.90\textwidth]{0.8\textheight}
{figures/xai_energy_svr_waveform.pdf}
{Linear-SVR importance along the energy-channel waveform.}

\vspace{-0.25em}


\mainmsg{Most of the useful timing information is concentrated
close to the signal arrival}.


\end{frame}
% ---------------------------------------------------------
\begin{frame}{XAI: timing-channel information}
\small

\centering
\smartimage[0.9\linewidth]{0.8\textheight}{figures/xai_timing_svr_waveform.pdf}{Timing-channel waveform importance.}

\vspace{-0.25em}

\mainmsg{Useful regions remain more than 20 ns after signal arrival.}

\end{frame}

% ---------------------------------------------------------
% ---------------------------------------------------------
% ---------------------------------------------------------
\begin{frame}{Conclusions}
\small

\begin{itemize}
    \item Using the dedicated timing channels, Pico-TDC measurements achieve
    CTR compatible with the oscilloscope reference.

    \vspace{0.45em}

    \item ML corrects the waveform-dependent timing error,
    with the largest improvement observed for the energy channels.

    \vspace{0.45em}

    \item The shared antisymmetric model gives a translation-invariant correction,
    while a few threshold-crossing times retain much of the useful information.

    \vspace{0.45em}

    \item Timing-relevant information can persist after signal arrival, potentially suggesting alternative timing strategies and electronics designs.
\end{itemize}

\vfill

\mainmsg{Machine learning can improve timing and reveal where useful information is encoded in the detector signal.}
\end{frame}
% ---------------------------------------------------------
\begin{frame}{References and source material}
\footnotesize

\mainmsg{Thank you for your attention.}

\begin{block}{References}
\begin{itemize}
  \item E. Berg and S. R. Cherry, ``Using convolutional neural networks to estimate time-of-flight from PET detector waveforms,'' \textit{Physics in Medicine \& Biology}, 2018.
  \item X. Feng, A. Muhashi, Y. Onishi, R. Ota, and H. Liu, ``Transformer-CNN hybrid network for improving PET time of flight prediction,'' \textit{Physics in Medicine \& Biology}, 2024.
  \item S. Naunheim et al., ``Explicit machine-learning corrections and spatial generalization for PET detector timing,'' \textit{Frontiers in Physics}, 2025.
  \item L. Y. Chang et al., ``A Time-Walk Correction Method for PET Detectors Based on Leading Edge Discriminators,'' \textit{IEEE Transactions on Radiation and Plasma Medical Sciences}, 2017.
\end{itemize}
\end{block}

\begin{block}{Repository}
\repo
\end{block}

\end{frame}



\end{document}
